\section*{Capítulo \ref{ch:b1:integration} --- Integração num Segmento}

\begin{solution}{exo:b1:integration:1}
$\dfrac{1}{x^2-4} = \dfrac{1/4}{x - 2} - \dfrac{1/4}{x+2}$
(multiplicação-avaliação), de modo que
\[
\int_0^1 \frac{\dd x}{x^2 - 4}
= \frac14\Bigl[\ln\abs{x-2} - \ln\abs{x+2}\Bigr]_0^1
= \frac14\Bigl(\ln\frac{1}{3} - \ln 1\Bigr) = -\frac{\ln 3}{4}.
\]

Por partes ($u' = 1$, $v = \ln t$): $\int_1^{\eu} \ln t\,\dd t =
\bigl[t\ln t\bigr]_1^{\eu} - \int_1^{\eu} \dd t = \eu - (\eu - 1) =
1$.

Por partes ($u' = \sin t$, $v = t$): $\int_0^\pi t\sin t\,\dd t =
\bigl[-t\cos t\bigr]_0^\pi + \int_0^\pi \cos t\,\dd t = \pi + 0 =
\pi$.
\end{solution}

\begin{solution}{exo:b1:integration:2}
Substituição $u = t^2 + 1$, $\dd u = 2t\,\dd t$:
\[
\int_0^1 \frac{t\,\dd t}{(t^2+1)^2}
= \frac12 \int_1^2 \frac{\dd u}{u^2}
= \frac12\Bigl[-\frac1u\Bigr]_1^2 = \frac14 .
\]
Com $u = \sin t$: $\int_0^{\pi/2} \cos^3 t\,\dd t = \int_0^{\pi/2}
(1 - \sin^2 t)\cos t\,\dd t = \bigl[\sin t - \frac{\sin^3
t}{3}\bigr]_0^{\pi/2} = 1 - \frac13 = \frac23$.
\end{solution}

\begin{solution}{exo:b1:integration:3}
$u_n = \dfrac1n \sum_{k=1}^{n} \dfrac{1}{1 + (k/n)^2}$: uma soma de Riemann
de $x \mapsto \frac{1}{1+x^2}$ em $\intcc{0}{1}$, de modo que $u_n \to
\int_0^1 \frac{\dd x}{1+x^2} = \arctan 1 = \dfrac\pi4$.

$\ln v_n = \dfrac1n \sum_{k=1}^{n} \ln\dfrac{n+k}{n} = \dfrac 1n
\sum_{k=1}^n \ln\Bigl(1 + \dfrac kn\Bigr) \to \int_0^1 \ln(1+x)\,\dd
x = \bigl[(1+x)\ln(1+x) - x\bigr]_0^1 = 2\ln 2 - 1$. Portanto, $v_n \to
\eu^{2\ln 2 - 1} = \dfrac 4\eu$. (Verificação da identificação:
$\frac{(2n)!}{n!\,n^n} = \prod_{k=1}^{n} \frac{n+k}{n}$.)
\end{solution}

\begin{solution}{exo:b1:integration:4}
O limite é $0$. Sejam $M = \sup \abs f$ e $\varepsilon \in
\intoo{0}{1}$. Corte em $1 - \varepsilon$:
\[
\Bigl| \int_0^1 x^n f \Bigr|
\leq \int_0^{1 - \varepsilon} x^n \abs f + \int_{1-\varepsilon}^1
x^n \abs f
\leq M\,(1-\varepsilon)^n + M\varepsilon .
\]
Como $(1 - \varepsilon)^n \to 0$ (\cref{exo:b1:seq:3}), o limite superior
do lado esquerdo é $\leq M\varepsilon$ para todo $\varepsilon$:
a \omterm{thm:b1:integration:def}{integral} tende a $0$.
\end{solution}

\begin{solution}{exo:b1:integration:5}
$Q(\lambda) = \int_a^b (f + \lambda g)^2 = \int f^2 + 2\lambda \int
fg + \lambda^2 \int g^2 \geq 0$ para todo $\lambda$. Se $\int g^2 =
0$, então $g = 0$ (positividade estrita,
\cref{thm:b1:integration:props}~(4)), e a desigualdade é $0 \leq
0$. Caso contrário, $Q$ é um genuíno \omterm{def:b1:poly:def}{polinômio} do segundo grau,
sempre $\geq 0$: o seu discriminante é $\leq 0$, isto é, $\bigl(\int fg\bigr)^2 \leq \int
f^2 \int g^2$.

Igualdade se, e somente se, o discriminante se anula, se, e somente se,
$Q(\lambda_0) = 0$ para algum $\lambda_0$, isto é, $\int (f + \lambda_0 g)^2 = 0$, ou seja (de novo por positividade
estrita), $f = -\lambda_0 g$: a igualdade vale exatamente quando
$f$ e $g$ são proporcionais.
\end{solution}

\begin{solution}{exo:b1:integration:6}
Seja $\Phi(a) = \int_a^{a+T} f$. Pelo teorema fundamental
(\cref{thm:b1:integration:ftc}), $\Phi$ é \omterm{def:b1:derivative:def}{derivável} com
$\Phi'(a) = f(a + T) - f(a) = 0$: constante.

Para $x > 0$, escreva $x = nT + r$, $0 \leq r < T$ ($n = \lfloor x/T
\rfloor$). Chasles:
\[
\int_0^x f = n \int_0^T f + \int_{nT}^{nT + r} f,
\qquad
\Bigl| \int_{nT}^{nT+r} f \Bigr| \leq T \sup_{\intcc{0}{T}} \abs f
= C .
\]
Então $\frac 1x \int_0^x f = \frac{nT}{x}\cdot\frac 1T \int_0^T f +
O\bigl(\frac 1x\bigr)$, e $\frac{nT}{x} \to 1$: o limite é
$\frac1T \int_0^T f$.
\end{solution}

\begin{solution}{exo:b1:integration:7}
Seja $g(x) = f(x) - x$: \omterm{def:b1:continuity:continuous}{contínua}, com
\[
\int_0^1 g = \int_0^1 f - \frac12 = 0 .
\]
Se $g$ nunca se anulasse, o teorema do valor intermediário forçaria um
sinal constante (uma função \omterm{def:b1:continuity:continuous}{contínua} num \omterm{prop:b1:reals:intervals}{intervalo} que assume os dois
sinais se anula); digamos $g > 0$. Então, pela positividade estrita
(\cref{thm:b1:integration:props}~(4) aplicada a $g > 0$, dando
$\int g > 0$): contradição com $\int g = 0$. Logo, $g(c) = 0$ para
algum $c$: $f(c) = c$.
\end{solution}

\begin{solution}{exo:b1:integration:8}
\begin{enumerate}
  \item Por partes, com $u' = \sin t$, $v = \sin^{n-1} t$:
        \[
        W_n = \bigl[-\cos t\sin^{n-1}t\bigr]_0^{\pi/2}
        + (n-1)\int_0^{\pi/2} \cos^2 t\,\sin^{n-2} t\,\dd t
        = (n-1)(W_{n-2} - W_n),
        \]
        de modo que $nW_n = (n-1)W_{n-2}$. De $W_0 = \frac\pi2$, $W_1 = 1$:
        \[
        W_{2p} = \frac{(2p-1)(2p-3)\cdots 1}{(2p)(2p-2)\cdots 2}\,
        \frac{\pi}{2} = \frac{(2p)!}{4^p (p!)^2}\,\frac\pi2,
        \qquad
        W_{2p+1} = \frac{(2p)(2p-2)\cdots 2}{(2p+1)(2p-1)\cdots 3}
        = \frac{4^p (p!)^2}{(2p+1)!} .
        \]
  \item Em $\intoo{0}{\frac\pi2}$, $0 < \sin t < 1$, de modo que $\sin^{n+1} <
        \sin^n$ e $(W_n)$ é (estritamente) decrescente e positiva.
        Comprimindo com a recorrência:
        \[
        \frac{n}{n+1} = \frac{W_{n+1}}{W_{n-1}}
        \leq \frac{W_{n+1}}{W_n} \leq 1
        \quad\implies\quad \frac{W_{n+1}}{W_n} \to 1 .
        \]
        Invariante: $a_n = (n+1)W_{n+1}W_n$ satisfaz $a_n = a_{n-1}$
        pela recorrência $(n+1)W_{n+1} = nW_{n-1}$, de modo que $a_n = a_0 =
        1 \cdot W_1 W_0 = \frac\pi2$. Então
        \[
        n W_n^2 \sim (n+1) W_{n+1} W_n = \frac\pi2
        \quad\implies\quad
        W_n \sim \sqrt{\frac{\pi}{2n}} .
        \]
\end{enumerate}
\end{solution}

\begin{solution}{exo:b1:integration:9}
\begin{enumerate}
  \item Em $\intoo{0}{\pi}$: $x > 0$, $a - bx = b(\frac ab - x) =
        b(\pi - x) > 0$ e $\sin x > 0$, de modo que o integrando é $> 0$
        e $I_n > 0$ (positividade estrita). Estimativa: em
        $\intcc{0}{\pi}$, $x \leq \pi$ e $a - bx \leq a$, de modo que $P
        \leq \frac{\pi^n a^n}{n!}$ e $I_n \leq
        \pi\,\frac{(\pi a)^n}{n!}$, que tende a $0$ (o
        fatorial vence o termo geométrico: é o termo geral
        da série exponencial convergente, cf.\
        \cref{ex:b1:seq:e}); em particular, $I_n < 1$ para $n$ grande.
  \item Expanda $x^n(a - bx)^n = \sum_{j=0}^{n} \binom nj a^{n-j}
        (-b)^j x^{n+j}$: assim, $P = \frac{1}{n!}\sum_j c_j x^{n+j}$
        com $c_j$ inteiros. Então $P^{(k)}(0) = 0$ para $k < n$
        (valoração) e, para $n \leq k \leq 2n$, $P^{(k)}(0) =
        \frac{k!}{n!} c_{k-n}$, um inteiro, pois $n! \mid k!$.
        Além disso, $P(\pi - x) = P(x)$ (substitua: $\pi - x$ troca
        os fatores, usando $a - b(\pi - x) = bx$), de modo que $P^{(k)}(\pi)
        = \pm P^{(k)}(0)$: também inteiros.
  \item Com $Q = P - P'' + P^{(4)} - \dots$ (finita: $P$ tem grau
        $2n$): $Q + Q'' = P$, e
        \[
        \bigl(Q'\sin x - Q\cos x\bigr)' = (Q + Q'')\sin x = P\sin x .
        \]
        Portanto, $I_n = \bigl[Q'\sin x - Q\cos x\bigr]_0^{\pi} = Q(\pi)
        + Q(0)$, uma soma de valores $P^{(2k)}$ em $0$ e em $\pi$: um
        inteiro, por (2).
  \item Para $n$ grande, $I_n$ é um inteiro com $0 < I_n < 1$:
        impossível. A suposição $\pi = \frac ab$ falha: $\pi$ é
        irracional.
\end{enumerate}
\end{solution}

\begin{solution}{exo:b1:integration:10}
Caso $C^1$: por partes,
\[
\int_a^b f(t)\sin\lambda t\,\dd t
= \Bigl[-f(t)\frac{\cos\lambda t}{\lambda}\Bigr]_a^b
+ \frac{1}{\lambda}\int_a^b f'(t)\cos\lambda t\,\dd t,
\]
limitado em valor absoluto por $\frac{2\sup\abs f + (b - a)\sup\abs{f'}}
{\lambda} \to 0$.

Caso \omterm{def:b1:continuity:continuous}{contínuo}: seja $\varepsilon > 0$ e tome uma \omterm{def:b1:integration:step}{função escada}
$\varphi$ com $\abs{f - \varphi} \leq \varepsilon$
(o \cref{thm:b1:integration:approx} fornece $\varphi \leq f \leq
\psi$ com folga $\leq\varepsilon$; tome $\varphi$). Então
\[
\Bigl|\int_a^b f\sin\lambda t\Bigr|
\leq \int_a^b \abs{f - \varphi}
+ \Bigl|\int_a^b \varphi \sin\lambda t\Bigr|
\leq (b-a)\varepsilon
+ \sum_i \abs{c_i}\,\Bigl|\int_{x_{i-1}}^{x_i} \sin\lambda t\,\dd
t\Bigr| ,
\]
e cada $\bigl|\int \sin \lambda t\bigr| = \bigl|\frac{\cos\lambda
x_{i-1} - \cos\lambda x_i}{\lambda}\bigr| \leq \frac{2}{\lambda}$:
o segundo termo tende a $0$. Portanto, o limite superior é $\leq
(b-a)\varepsilon$ para todo $\varepsilon$: o limite é $0$.
\end{solution}

\begin{solution}{exo:b1:integration:11}
Sejam $m = \min f$ e $M = \max f$, atingidos pelo teorema de
Weierstrass. Como $g \geq 0$: $m\,g \leq fg \leq M\,g$, de modo que, por
monotonicidade,
\[
m \int_a^b g \;\leq\; \int_a^b fg \;\leq\; M \int_a^b g .
\]
Se $\int_a^b g = 0$: a positividade estrita
(\cref{thm:b1:integration:props}~(4)) força $g \equiv 0$, os dois
lados se anulam, e qualquer $c$ serve. Caso contrário, $t = \frac{\int
fg}{\int g}$ está em $\intcc{m}{M} = f(\intcc{a}{b})$
(\cref{thm:b1:continuity:evt,thm:b1:continuity:ivt}), de modo que $t =
f(c)$ para algum $c$.

O sinal importa: em $\intcc{-1}{1}$, com $f(t) = g(t) = t$: $\int
fg = \int_{-1}^1 t^2 = \frac23$, ao passo que $f(c)\int_{-1}^1 t\,\dd t
= 0$ para todo $c$.
\end{solution}

\begin{solution}{exo:b1:integration:12}
Se $f \equiv 0$, a afirmação é vazia (todo ponto é zero). Suponha, então,
$f \not\equiv 0$ e que ela tenha no máximo $n$ zeros distintos
em $\intoo{a}{b}$; sejam $z_1 < \dots < z_m$ ($m \leq n$) os
zeros em que $f$ \emph{muda de sinal} (possivelmente nenhum). Ponha
$P(t) = \prod_{i=1}^{m}(t - z_i)$ (produto vazio $= 1$), de
grau $m \leq n$. Em cada subintervalo cortado pelos $z_i$, tanto
$f$ quanto $P$ têm sinal constante, e ambos trocam de sinal ao cruzar
algum $z_i$: o produto $fP$ tem um único sinal constante em todo
$\intoo{a}{b}$. Sendo \omterm{def:b1:continuity:continuous}{contínuo}, não identicamente nulo e de
sinal constante, ele tem $\bigl|\int_a^b fP\bigr| > 0$ (positividade
estrita aplicada a $\abs{fP}$). Mas $\int f P$ é uma combinação
linear dos momentos $\int f\,t^k$, $k \leq n$, todos nulos:
contradição. Portanto, $f$ tem ao menos $n + 1$ zeros distintos em
$\intoo{a}{b}$.
\end{solution}

\begin{solution}{pb:b1:integration:1}
\textbf{1.} Seja $N = \lceil 2c \rceil$. Para $n \geq N$:
$\frac{c^{n+1}/(n+1)!}{c^n/n!} = \frac{c}{n+1} \leq \frac12$, de modo que
$0 < \frac{c^n}{n!} \leq \frac{c^N}{N!}\,2^{-(n - N)} \to 0$:
confronto.

\textbf{2.} Fixe $k$; indução em $m$. Para $m = 0$: $\int_0^1
x^k = \frac{1}{k+1} = \frac{k!\,0!}{(k+1)!}$. Passo, por partes
($u = (1-x)^m$, $v' = x^k$):
\[
\int_0^1 x^k (1-x)^m \dd x
= \Bigl[\frac{x^{k+1}}{k+1}(1-x)^m\Bigr]_0^1
+ \frac{m}{k+1}\int_0^1 x^{k+1}(1-x)^{m-1}\dd x
= \frac{m}{k+1}\cdot\frac{(k+1)!\,(m-1)!}{(k+m+1)!} ,
\]
que é $\frac{k!\,m!}{(k+m+1)!}$.

\textbf{3.} $k = m = n$: $\int_0^1 (x(1-x))^n =
\frac{(n!)^2}{(2n+1)!} = \frac{1}{(2n+1)\binom{2n}{n}}$. Como
$x(1-x) \leq \frac14$ em $\intcc{0}{1}$, a \omterm{thm:b1:integration:def}{integral} é $\leq
4^{-n}$, donde $\binom{2n}{n} \geq \frac{4^n}{2n+1}$ ---
coerente com $\binom{2n}{n}^{1/n} \to 4$
(\cref{pb:b1:seq:1}).

\textbf{4.} O integrando é \omterm{def:b1:continuity:continuous}{contínuo}, $\geq 0$ e positivo
em $\intoo{0}{1}$, logo não identicamente nulo: a sua \omterm{thm:b1:integration:def}{integral} é
$> 0$ (\cref{thm:b1:integration:props}~(4)). Estimativa superior:
$(x(1-x))^n \leq 4^{-n}$ e $g \leq \sup\abs g$, e depois
monotonicidade.

\textbf{5.} $A_0 = \eu - 1$; $A_1 = [x\eu^x]_0^1 - \int_0^1
\eu^x = \eu - (\eu - 1) = 1$. Por partes: $A_n = [x^n \eu^x]_0^1 -
n\int_0^1 x^{n-1}\eu^x = \eu - n A_{n-1}$. Estimativas: o integrando
é positivo, de modo que $A_n > 0$; e $\eu^x \leq \eu$ dá $A_n \leq
\eu\int_0^1 x^n = \frac{\eu}{n+1}$.

\textbf{6.} $A_0 = -1 + 1\cdot\eu$. Se $A_{n-1} = \alpha_{n-1} +
\beta_{n-1}\eu$ com entradas inteiras, então
\[
A_n = \eu - n\alpha_{n-1} - n\beta_{n-1}\eu
= \underbrace{(-n\,\alpha_{n-1})}_{\alpha_n}
+ \underbrace{(1 - n\,\beta_{n-1})}_{\beta_n}\,\eu ,
\]
ambos inteiros.

\textbf{7.} Se $\eu = \frac pq$: $q A_n = q\alpha_n + p\beta_n
\in \Z$, e $0 < qA_n \leq \frac{q\eu}{n+1} < 1$ para $n$ grande:
um inteiro estritamente entre $0$ e $1$ --- impossível. Logo, $\eu
\notin \Q$. No \cref{exo:b1:seq:9}, o inteiro preso era
$q!\,\frac pq - q!\,a_q$; aqui é $qA_n$: mesma armadilha,
combustível \omterm{thm:b1:integration:def}{integral}.

\textbf{8.} $n = 0$: $R_0 = \int_0^1 \eu^t = \eu - 1$, de modo que $\eu =
1 + R_0$. Por partes ($u = \eu^t$, $v = -\frac{(1-t)^{n+1}}{n+1}$):
\[
R_n = \frac{1}{n!}\Bigl(\Bigl[-\frac{(1-t)^{n+1}}{n+1}
\eu^t\Bigr]_0^1 + \frac{1}{n+1}\int_0^1 (1-t)^{n+1}\eu^t\Bigr)
= \frac{1}{(n+1)!} + R_{n+1} ,
\]
de modo que a fórmula se propaga de $n$ a $n + 1$. Estimativas: $1 \leq
\eu^t \leq \eu$ em $\intcc{0}{1}$ e $\int_0^1 (1-t)^n =
\frac{1}{n+1}$ dão $\frac{1}{(n+1)!} \leq R_n \leq
\frac{\eu}{(n+1)!}$.

\textbf{9.} $\sum_{k=0}^{9} \frac{1}{k!} = \frac{986410}{362880}
= 2.71828152\dots$, e
\[
\frac{1}{10!} = 2.76\cdot10^{-7} \leq R_9 \leq
\frac{2.75}{10!} = 7.58\cdot10^{-7} ,
\]
de modo que $2.7182818 \leq \eu \leq 2.7182823$: com demonstração, $\eu =
2.718282$ com seis casas decimais (valor verdadeiro $2.7182818\dots$).

\textbf{10.} $f(1 - x) = \frac{(1-x)^n x^n}{n!} = f(x)$. Em
$\intoo{0}{1}$: $0 < x(1-x) \leq \frac14$, de modo que $0 < f \leq
\frac{1}{4^n\,n!}$.

\textbf{11.} $x^n(1-x)^n = \sum_{j=0}^{n} (-1)^j\binom
nj\,x^{n+j}$, de modo que $f = \frac{1}{n!}\sum_j c_j\,x^{n+j}$ com
$c_j \in \Z$. Portanto, $f^{(k)}(0) = 0$ para $k < n$ ou $k > 2n$,
e, para $n \leq k \leq 2n$: $f^{(k)}(0) = \frac{k!}{n!}\,
c_{k-n}$, um inteiro, pois $n! \mid k!$. A simetria dá
$f^{(k)}(1) = (-1)^k f^{(k)}(0) \in \Z$.

\textbf{12.} $b^n \pi^{2n-2k} = b^n\bigl(\frac
ab\bigr)^{\,n-k} = a^{\,n-k}\,b^{\,k} \in \Z$, de modo que $G(0) =
\sum_k (-1)^k a^{n-k} b^k f^{(2k)}(0)$ e, do mesmo modo, $G(1)$ são
inteiros pela questão 11.

\textbf{13.} Em $\pi^2 G + G''$, o termo $k$ de $\pi^2 G$
carrega $\pi^{2n-2k+2}f^{(2k)}$, e o termo $j = k - 1$ de
$G''$ carrega $(-1)^{k-1}\pi^{2n-2k+2}f^{(2k)}$: tudo se
cancela, exceto $k = 0$ na primeira soma e $j = n$ na
segunda, isto é,
\[
G'' + \pi^2 G = b^n\bigl(\pi^{2n+2} f + (-1)^n f^{(2n+2)}\bigr)
= b^n \pi^{2n+2} f = \pi^2 a^n f
\]
($f$ tem grau $2n$, de modo que $f^{(2n+2)} = 0$; e $b^n\pi^{2n} =
a^n$). Então
\[
\bigl(G'\sin\pi x - \pi G\cos\pi x\bigr)'
= (G'' + \pi^2 G)\sin \pi x = \pi^2 a^n f\sin\pi x .
\]

\textbf{14.} Integrando em $\intcc{0}{1}$:
\[
\pi^2 a^n \int_0^1 f\sin(\pi x)\,\dd x
= \bigl[G'\sin\pi x - \pi G\cos\pi x\bigr]_0^1
= \pi\bigl(G(1) + G(0)\bigr) ,
\]
de modo que $\pi a^n \int_0^1 f\sin\pi x = G(0) + G(1) \in \Z$. Em
$\intoo{0}{1}$, $f > 0$ e $\sin\pi x > 0$: o lado esquerdo é
positivo, logo $G(0) + G(1) \geq 1$; e $\sin \leq 1$, com a
questão 10, o limita por $\frac{\pi a^n}{4^n\,n!}$.

\textbf{15.} Pela questão 1 (com $c = \frac a4$),
$\frac{\pi a^n}{4^n n!} \to 0$: para $n$ grande, ele é $< 1$,
contradizendo $G(0) + G(1) \geq 1$. Assim, fração $\frac ab$ alguma
é igual a $\pi^2$: o teorema de Legendre. Se $\pi$ fosse racional,
$\pi^2$ também seria: logo, $\pi \notin \Q$ --- e a implicação
só corre nesse sentido ($\sqrt2$ é irracional com quadrado
racional), e é por isso que $\pi^2 \notin \Q$ é estritamente mais forte
que o \cref{exo:b1:integration:9}.

\textbf{16.} Integralidade: questões 11--12 (\omterm{def:b1:derivative:def}{derivadas} nas
extremidades); pequenez: questões 10 e 1; ponte: questões
13--14 (a dupla \omterm{thm:b1:integration:parts}{integração por partes} telescopada). Sem
$\frac{1}{n!}$, os dados nas extremidades continuam inteiros (com ainda
mais facilidade), mas a estimativa se torna $\frac{\pi a^n}{4^n}$, que
tende a $0$ apenas quando $a < 4$ --- e todo candidato tem $a =
b\,\pi^2 > 9$. O fatorial é exatamente o que ultrapassa o
crescimento geométrico $a^n$: sem fatorial, sem teorema.

\textbf{17.} Para $a \leq 10$, o inteiro $G(0) + G(1)$ é
positivo e no máximo $\pi\,(10/4)^n/n!$. Em $n = 7$:
$2.5^7 = 610.35\dots$, de modo que a estimativa é $\frac{\pi \times
610.35}{5040} \approx 0.38 < 1$ (em $n = 6$ ela ainda vale
$1{,}07$): a contradição aparece no sétimo estágio,
explicitamente.

\textbf{18.} Duas \omterm{thm:b1:integration:parts}{integrações por partes}:
\[
\int_0^1 x(1-x)\sin\pi x\,\dd x
= \frac1\pi\int_0^1 (1 - 2x)\cos\pi x\,\dd x
= \frac{2}{\pi^2}\int_0^1 \sin\pi x\,\dd x
= \frac{2}{\pi^2}\cdot\frac{2}{\pi} = \frac{4}{\pi^3}
\]
(os termos de \omterm{def:b1:topology:closure}{fronteira} se anulam: $x(1-x)$ em $0, 1$, e $\sin\pi x$
em $0, 1$). Simbolicamente, o telescópio $n = 1$ (sem suposição
sobre $\pi$) diz $\pi^3\int_0^1 f_1\sin\pi x = -(f_1''(0) +
f_1''(1))$ com $f_1 = x(1 - x)$, $f_1'' = -2$: lado direito $4$
--- os dois cálculos concordam.

\textbf{19.} $\eu^x$ resolve $y' = y$ e $\sin\pi x$ resolve
$y'' = -\pi^2 y$: equações lineares com coeficientes constantes,
de modo que a \omterm{thm:b1:integration:parts}{integração por partes} faz o núcleo circular de volta a si
mesmo e mantém todos os dados de \omterm{def:b1:topology:closure}{fronteira} dentro de $\Z + \Z\eu$ (resp.\
\omterm{def:b1:poly:def}{polinômios} inteiros em $\pi^2$). Essa propriedade de fechamento é o que
a máquina precisa. Refinado com núcleos adaptados a vários pontos ao
mesmo tempo, o mesmo mecanismo produz o teorema de Hermite ($\eu$
transcendente, 1873) e o de Lindemann ($\pi$ transcendente,
1882) --- além deste volume.

\textbf{20.} Divisão polinomial (ou multiplique de volta e confira):
\[
x^4(1-x)^4 = (x^6 - 4x^5 + 5x^4 - 4x^2 + 4)(1 + x^2) - 4 .
\]
Integrando a identidade exibida, dividida por $1 + x^2$:
\[
\int_0^1 \frac{x^4(1-x)^4}{1+x^2}\dd x
= \Bigl(\frac17 - \frac46 + 1 - \frac43 + 4\Bigr)
- 4\arctan 1
= \frac{22}{7} - \pi ,
\]
usando $\frac17 - \frac23 + 1 - \frac43 + 4 = 3 + \frac17$ e
$\arctan 1 = \frac\pi4$.

\textbf{21.} O integrando é \omterm{def:b1:continuity:continuous}{contínuo} e positivo em
$\intoo{0}{1}$: a \omterm{thm:b1:integration:def}{integral} é $> 0$, de modo que $\pi < \frac{22}{7}$.
Além disso, $\frac12 \leq \frac{1}{1+x^2} \leq 1$ em
$\intcc{0}{1}$ e $\int_0^1 (x(1-x))^4 = \frac{(4!)^2}{9!} =
\frac{1}{630}$ (questão 3):
\[
\frac{1}{1260} \leq \frac{22}{7} - \pi \leq \frac{1}{630}
\quad\Longrightarrow\quad
3.14126 < \frac{22}{7} - \frac{1}{630} \leq \pi \leq
\frac{22}{7} - \frac{1}{1260} < 3.14207 .
\]

\textbf{22.} \omterm{def:b1:complex:field}{Módulo} $x^2 + 1$: $x^2 \equiv -1$, de modo que $x^{4m} =
(x^2)^{2m} \equiv 1$ e $(1 - x)^2 = 1 - 2x + x^2 \equiv -2x$,
donde $(1-x)^{4m} \equiv (-2x)^{2m} = 4^m (x^2)^m \equiv
(-4)^m$: o resto é a constante $(-4)^m$, e o
quociente $Q_m$ tem coeficientes inteiros (divisão por um \omterm{def:b1:poly:def}{polinômio}
inteiro \omterm{def:b1:poly:def}{mônico}). Dividindo a identidade por $1 + x^2$ e
integrando:
\[
J_m := \int_0^1 \frac{(x(1-x))^{4m}}{1+x^2}\dd x
= s_m + (-4)^m\,\frac{\pi}{4},
\qquad s_m = \int_0^1 Q_m \in \Q .
\]
Resolvendo para $\pi$: com $r_m = (-1)^{m+1}\,4^{\,1-m} s_m \in
\Q$, $\;\abs{\pi - r_m} = 4^{\,1-m} J_m \leq
4^{\,1-m}\cdot4^{-4m} = 4^{\,1-5m}$. Para $m = 1$: $s_1 =
\frac{22}{7}$, $r_1 = \frac{22}{7}$, estimativa $4^{-4} =
\frac{1}{256}$ --- de novo as questões 20--21.

\textbf{23.} $\bigl|\pi - \frac{22}{7}\bigr| = \frac{22}{7} -
\pi \approx 1.26\cdot10^{-3}$, dezesseis vezes melhor que o
padrão de ordem $2$, $\frac{1}{7^2} \approx 2.0\cdot10^{-2}$,
garantido por Dirichlet (\cref{pb:b1:derivative:1}, questão 4);
e $\bigl|\pi - \frac{355}{113}\bigr| \approx 2.7\cdot10^{-7}$
supera $\frac{1}{113^2} \approx 7.8\cdot10^{-5}$ por um fator
$\approx 300$. Nenhuma contradição com o que foi demonstrado: as
desigualdades de Liouville \emph{limitam inferiormente} os erros de
aproximação apenas para \omterm{pb:b1:logic:1}{números algébricos}, e nenhuma estimativa
dessas para $\pi$ está disponível neste nível --- $\pi$ está livre
para ser aproximado espetacularmente bem.

\textbf{24.} Lema: suponha $x = \frac pq$ e $0 < \abs{a_n +
b_n x} \to 0$. Então $\abs{a_n + b_n x} = \frac{\abs{q a_n + p
b_n}}{q}$, com $q a_n + p b_n$ um inteiro não nulo (não nulo,
porque o valor absoluto é $> 0$): assim, $\abs{a_n + b_n x}
\geq \frac1q$ para todo $n$, contradizendo a convergência a
$0$. Instâncias: questão 7 ($x = \eu$, $a_n = \alpha_n$, $b_n =
\beta_n$); \cref{exo:b1:seq:9} ($x = \eu$ de novo, com $a_q =
-q!\sum_{k \leq q}\frac{1}{k!}$, $b_q = q!$); e o
\cref{pb:b1:derivative:1}, questão 1, é a sua forma geométrica.
Na Parte III e no \cref{exo:b1:integration:9}, a armadilha roda
\emph{dentro} da contradição: supor a racionalidade converte
uma expressão num inteiro, que a análise em seguida comprime
em $\intoo{0}{1}$ --- o mesmo princípio, transposto.

\textbf{25.} (i) A integralidade vive no cálculo nas extremidades de
\omterm{def:b1:poly:def}{polinômios} (questões 6, 11--12), a pequenez nas estimativas que a
monotonicidade do $\sup$ dá (questões 4 e 10), e a ponte na
\omterm{thm:b1:integration:parts}{integração por partes} (questões 8, 13--14) --- os três são
teoremas deste capítulo. (ii) A \omterm{thm:b1:integration:def}{integral} fornece o que o
teorema do valor médio não podia: uma identidade \emph{exata} entre
o objeto analítico e os dados aritméticos (igualdade, e não apenas
uma desigualdade com um $c$ desconhecido), e é por isso que a máquina
alcança $\pi^2$, ao passo que o \cref{pb:b1:derivative:1} alcançou apenas
expoentes de aproximação. (iii) Extraídos: $\eu \notin \Q$,
$\pi^2 \notin \Q$ (logo, $\pi \notin \Q$), $\eu = 2.718282$
certificado, $\frac{22}{7} - \frac{1}{630} \leq \pi \leq
\frac{22}{7} - \frac{1}{1260}$, e $\binom{2n}{n} \geq
\frac{4^n}{2n+1}$. (iv) \omterm{def:b1:topology:closure}{Fronteira}: Hermite e Lindemann empurram
a mesma máquina até a transcendência; e Apéry (1979) rodou a
armadilha do inteiro em $\zeta(3)$ --- a máquina ainda está produzindo
matemática do século XX.
\end{solution}
